\begin{abstract}
\noindent
Prediction intervals for inflation inform decisions in commodity-dependent economies, but it is not known whether adaptive conformal methods that perform well on euro-area and US data keep their nominal coverage there. We build a multi-horizon (1, 3, 6 and 12 months) calibration benchmark of split, rolling-window, adaptive (ACI, FACI), conformal PID and regime-aware (DRACP) intervals at nominal level 0.90, using monthly IMF consumer price indices for 13 design economies, a first hold-out panel of 7 economies and, in an extension, a second hold-out panel of 10 economies, with LightGBM, ARIMA, a naive forecaster and the zero-shot foundation model Chronos as base forecasters. The primary outcome is the calibration error (CE), the mean over countries of $|\text{coverage}-0.90|$. Contrasts are tested with a moving-block bootstrap and Holm correction, equivalence of DRACP and ACI with two one-sided tests and a margin of 2 percentage points (pp), and decision rules R1--R5 were pre-specified in a version-controlled protocol before the hold-out run (a self-attested commit history; the rules were written after the exploratory results were known).

Adaptive methods had a lower calibration error than a calibration performed once and, in a post hoc comparison restricted to ARIMA, than the native intervals of that model; the direction is consistent across samples. The pre-specified rule R1 (ACI against split) was met only at a 96-month warm-up, a sensitivity run that is not promoted, and not at the primary 48-month warm-up (ACI minus split at $h=3$: $-3.0$~pp, $p=0.054$; at $h=6$: $-2.8$~pp, $p=0.092$). Whether ACI improves on a rolling window depends on the panel: it was not shown on the design panel and the first hold-out panel (minimum detectable difference undefined, inconclusive), and it was met under the pre-specified rule in part of the cells of the second hold-out panel at warm-up 48; the reason is not known and the panels are not pooled. DRACP against ACI is inconclusive, and because the 2~pp margin is about the size of ACI's own calibration error this verdict is weak by construction. No efficiency advantage in interval width or interval score is established. At the 12-month horizon pooled coverage is below nominal for every method on the design and first hold-out panels; in a diagnostic run that releases outcomes immediately, ACI coverage on the first hold-out panel was 2.4~pp higher (95\% interval $-0.1$ to $4.8$, which includes zero), below the pre-specified 3~pp threshold. On the 22 target months from 2025-01, a period that Chronos cannot be assumed to have seen in pretraining (its cutoff was not verified), its mean absolute error was lower than that of the naive baseline and, descriptively, at least as low as that of LightGBM and ARIMA. A language-model explanation layer was evaluated as an engineering extension that enters no main claim: one model, 468 prompts in two conditions, 936 answers, programmatic metrics. With the forecast payload in the prompt, the pre-specified metrics gave 90.6\% numerical faithfulness (424 of 468), 9.0\% omission of the interval and 23.3\% coverage overstatement. Two defects of these metrics, found by reading answers (an en-dash range read as a negative number, and negated guarantee words flagged), were corrected post hoc; the corrected version gave 100\% (468 of 468), 0\% and 3.0\%, the remaining overstatement flags being false positives, as judged on 60 answers by a single reader that was the model itself, not an independent human review. Without the payload no answer gave an interval, and 145 of the 156 answers that contain numbers have at least one number that matches nothing in the payload. This does not show that grounding works in general or that the model is faithful in general; Kazakh fluency and the human review were not done. The contributions are a benchmark with null outcomes specified in advance, a comparison of adaptive calibration with one-shot calibration and native intervals, a feedback-delay diagnostic, a panel-dependence finding that is reported rather than resolved, and an open reproducible pipeline.
\end{abstract}
